\section{Discussão crítica}
\label{sec:discussao}

Esta seção discute as limitações das decisões tomadas ao longo do trabalho e o quanto cada uma afeta os resultados.

\paragraph{Tratamento dos dados.} O grafo foi reduzido à lista de arestas, sem pesos e sem distinção entre os critérios de ligação (Seção~\ref{sec:tratamento}). Uma aresta entre imagens do mesmo local vale o mesmo que uma entre imagens com uma tag em comum, e um par que compartilha vários metadados vale o mesmo que um que compartilha um só. As medidas descrevem, portanto, a união de relações diferentes, e não permitem dizer qual critério produz a clusterização alta ou os hubs. Além disso, o arquivo só lista imagens com pelo menos uma aresta (Tabela~\ref{tab:basicas}), de modo que as imagens sem metadados em comum com nenhuma outra não fazem parte de $n$. As medidas de distância e a robustez usam só a maior componente, mas essa restrição pesa pouco, porque ela reúne quase todos os vértices e arestas (Tabela~\ref{tab:tratamento}).

\paragraph{Estimativas por amostragem.} Quase todas as medidas da Parte I são exatas (Tabela~\ref{tab:cobertura_parte1}). São estimados o comprimento médio da análise small-world, a partir de 1.000 fontes, e a métrica C da robustez, a partir de 100 fontes. No primeiro caso, o valor exato da Tabela~\ref{tab:distancias} está dentro do intervalo de confiança da Tabela~\ref{tab:small_world}, e a amostragem não afeta a conclusão. No segundo, a estimativa da maior componente original difere do valor exato só nos centésimos (Tabelas~\ref{tab:distancias} e~\ref{tab:robustez_comparacao}), bem menos que a diferença entre os cenários aleatório e direcionado. O cenário direcionado, porém, é uma única execução, sem medida de dispersão, e o desvio do cenário aleatório mistura a variação dos vértices removidos com a das fontes sorteadas. A visualização (Figura~\ref{fig:visualizacao}) é uma amostra parcial e só ilustra a estrutura local.

\paragraph{Política de pesos artificiais.} Como o grafo não tem pesos, os algoritmos de caminhos mínimos e de árvore geradora mínima usaram pesos inteiros sorteados uniformemente em $[1, 10]$ com semente fixa (Seção~\ref{sec:parte2}). Esses pesos não carregam informação sobre as imagens: as distâncias e o peso da árvore das Tabelas~\ref{tab:tempos_caminhos}, \ref{tab:tempos_floyd} e~\ref{tab:tempos_agm} são propriedades do sorteio, e não do grafo, e servem apenas para exercitar e validar os algoritmos. Também os tempos dependem dos pesos: o número de rodadas do Bellman-Ford é limitado pelo número de arestas dos caminhos mínimos, que muda com os pesos.

Os atributos de aresta que o SNAP distribui em \texttt{edgeFeaturesFlickr.tar.gz} não servem de peso. Cada par de imagens tem um vetor de atributos, com contagens de metadados em comum e indicadores, sem uma escala única, e reduzi-lo a um número exigiria uma combinação arbitrária. Os arquivos são separados pela coleção de origem das imagens (CLEF, MIR, PASCAL e NUS). Só parte das arestas de \texttt{flickrEdges.txt.gz} tem atributos, de modo que as demais ficariam sem peso, e a maioria dos pares com atributos não é aresta do grafo. Os pesos $w_{ij}$ do artigo de origem~\cite{mcauley2012image} também não servem: são potenciais aprendidos para cada categoria de rótulo, e não um peso por aresta.

\paragraph{Tempos medidos em uma única máquina.} Todas as medições foram feitas numa única máquina (Seção~\ref{sec:metodologia}). Os valores absolutos dependem do processador, da memória, do sistema operacional e das versões do Python e das bibliotecas, e mudariam em outra máquina. O intervalo de confiança cobre só a variação entre as repetições nessa máquina, e não entre máquinas. O que se transfere é a tendência de crescimento com o tamanho da entrada, comparada à complexidade teórica. Mesmo ela é afetada pela máquina: a hierarquia de memória é uma explicação compatível com a inclinação do Floyd-Warshall acima de 3 (Figura~\ref{fig:floyd_escala}), sem ter sido medida, e a interferência de outros processos é uma explicação possível para os desvios maiores de algumas entradas (Tabela~\ref{tab:tempos_buscas}). As comparações entre o networkx e o igraph (Tabela~\ref{tab:clusterizacao}) refletem sobretudo a linguagem de cada implementação, Python ou C, e não o algoritmo. Por fim, os subgrafos crescentes são aninhados e não crescem de forma uniforme em $m$, e cinco pontos por algoritmo dão uma tendência, não uma estimativa precisa da ordem de crescimento.

\paragraph{Interpretação de small-world.} As referências aleatórias preservam só $n$ e o grau médio, e não a distribuição de graus, fortemente assimétrica (Figura~\ref{fig:graus_loglog}). Uma referência que preservasse também os graus de cada vértice não foi avaliada, e os seus valores de $C_{\mathrm{rand}}$ e $L_{\mathrm{rand}}$ poderiam ser outros. Além disso, $C_{\mathrm{rand}} \approx \langle k \rangle/n$ diminui quando $n$ cresce com $\langle k \rangle$ fixo, de modo que $C/C_{\mathrm{rand}}$ e $\sigma$ tendem a ser grandes em qualquer grafo grande e esparso com clusterização apreciável. A magnitude de $\sigma$ (Tabela~\ref{tab:small_world}) reflete, em parte, o tamanho do grafo. A conclusão se apoia mais na comparação direta de $C$ e de $L$ com as referências. Por fim, a clusterização alta pode ser, em parte, consequência da própria regra de construção, que liga entre si as imagens que compartilham um metadado (Seção~\ref{sec:descricao}), e não uma propriedade que emerge da rede.

\paragraph{Interpretação da lei de potência.} As limitações do ajuste foram discutidas junto da resposta (Seção~\ref{sec:parte3}): o trecho ajustado cobre menos de uma década de $k$ e uma pequena fração dos vértices (Tabela~\ref{tab:lei_potencia}), não há teste de hipótese nem comparação com distribuições alternativas, e o estimador de $\alpha$ é aproximado. A elevação local da distribuição perto de $k = 100$ (Figura~\ref{fig:graus_loglog}) mostra, além disso, que a distribuição tem estrutura que nenhuma lei de potência única descreve. O resultado não permite, portanto, classificar o grafo como livre de escala.

\paragraph{Robustez.} A análise usou uma única fração de remoção, 5\%, e, no ataque direcionado, o ranking de grau calculado uma única vez, antes de qualquer remoção (Seção~\ref{sec:parte3}). Um ataque que recalcula o grau após cada remoção e ataques por outras centralidades, como a betweenness, evitada pelo custo, não foram avaliados. A conclusão de que a rede não é muito mais vulnerável ao ataque direcionado vale, portanto, para esse ataque e essa fração, pelo critério de $S$ e do comprimento médio. Pelas métricas B e D, o ataque fragmenta a periferia muito mais que a falha aleatória (Tabela~\ref{tab:robustez_comparacao}).
